\subsection{The Network Transmission Model}
\label{subsec:net}

Contemporary network theories of cultural taste and attitude formation and change (and by implication decay) suggest that the contemporary taste profile of an individual is largely a function of the tastes held by the other persons in her surrounding social network \citep{erickson82, erickson88}.  The key underlying assumption is that social ties serves as the ``conduits'' through which new cultural practices are transmitted and by way of which individuals learn these new cultural practices from their social contacts \citep{mark98a, carley91}.  This proposition is shared by both straightforward network accounts \citep[e.g.][]{cardon_granjon05, erickson96, kane04} and more ambitious ecological-network theories of cultural taste \citep{mark98a, mark03, mcpherson04}. \citet[323--324]{mark03}  for instance speaks of social ties to friends and family as the most important sources of cultural taste acquisition:  ``[p]eople's friends and family{\ldots}influence the leisure activities they participate in{\ldots}participation in various sports, visiting art museums, and gardening are{\ldots}leisure activities that people often lean from parents, siblings, and/or friends.''

Network theories of taste transmission implicitly presume a relatively stable ``social structure'' coupled to relatively unstable patterns of cultural taste, otherwise it would be hard to think how something that varies at a more rapid pace ``causes'' or ``affects'' something that is more stable than itself \citep{simon52}.  It bears mentioning that this relative stability assumption is shared by the entire family of network theories that conceive of social structure as the pipes through which contents, or information ``flows'' and diffuses from one actor to another as \citet{borgatti05} has recently noted.  This implicit (and crucial) assumption of the stability of social structure, when this term is defined---as it is usually done in network theory---as the ``concrete'' patterns of network relations that bind individuals to one another \citep{berkowitz82, berkowitz88, wellman83, wellman88, bearman97, breiger04}, can certainly be questioned in light of recent longitudinal research on constancy of network relations over time.  This research---see the 1997 and 2005 special issues of \textit{Social Networks} (volume 19, number 1 and volume 27 number 4) and \citet{burt00} for a thorough review)---shows that instead of being constant, network relations are extremely volatile, even when the time frame is constrained to relatively short periods of time, such as a year or less \citet{bidart_degenne05, suitor_etal97}. Thus, what is truly surprising about over-time network ``stability'' over time is ``how unstable intimacy is'' \citep[47]{wellman_etal97}.  

Accordingly, most longitudinal studies of personal networks reveal surprisingly high levels of instability in personal networks over time \citep{lang00}, with anywhere approximately about 30 to 60\% of the personal network experiencing turnover in about a year's time.  In his own research on bankers, \citet[11]{burt00} finds that ``three in four of the 22,709 colleague relations at risk of decay{\ldots}are gone next year.'' This means that old ties are constantly being deleted---what \citet{burt00} refers to as ``tie decay''---and new ones formed constantly throughout the life-course \citet{suitor_etal97, wellman_etal97}.  \citet[57]{suitor_keeton97}, report that more than half of the associates...named at T1 were not named again at either T2 or T3.''   \citet[14]{morgan_etal97} conclude that their results ``point to a considerable degree of instability in who is present in the network over time. In particular, the average overlap rate of 55\% means that just over half of those who were named in one interview or the other appeared in both. Similarly, interpreting the average phi value of 0.34 as a test-retest correlation indicates only a limited ability to predict who was in the network from one time to the next.''

These findings imply that for most individuals a recurrent process of new contact selection, extant relationship ``fine tuning'' and possible deletion of old social ties is constantly at work \citep{lazarsfeld_merton54}.  This represents an important---if ill understood---facet underlying the temporal evolution of a person's relational micro-environment.   This dynamic acquires renewed salience in light of the observation that recently formed network connections appear to exhibit what \citet{burt00}, drawing on \citet[148--150]{stinchcombe65}, refers to as a ``liability of newness.'' This means that dyadic network ties tend to display a high probability of dissolution at early stages of the relationship, with the probability of termination steeply declining as the relationship matures. This proposition implies that the ``tie decay function'' (that is the probability of a tie being dissolved) is non-linearly decreasing in relationship age, with a high probability for young relationships which decreases as these dyadic ties ripen into close friendships.  \citet{burt00, burt02} finds evidence consistent with this hypothesis for both strong ties and network bridges (weak ties), respectively.

It is true that network stability is much more substantial in the ``core'' personal network, with the bulk of the instability concentrated on the periphery\citep{bidart_lavenu05, suitor_etal97, wellman_etal97}.  However, most longitudinal network studies show that the stable ``core'' network tends to be composed of a small set of alters, primarily kin.  The periphery of the network includes most of the persons usually referred to as (non-kin) ``friends.''  Studies show that there is considerable turnover even among those ties usually considered ``close friends'' \citep{bidart_lavenu05}. It is possible to revise network theories of taste transmission to suggest that taste are mostly transmitted through the small, relatively stable core and not through the larger, unstable periphery.  However, most contemporary network theories of taste do not make this distinction, and assume that taste are as equally likely to be learned from family and friends \citep{mark03}.  It is clear from recent studies on personal communities that the ``core'' kin network while mobilized for many purposes, including emotional and financial support, is not used for companionship and sociability \citep{simmel49} in any meaningful sense.  Most research and theory on culture consumption communities shows that if tastes are exploited in interaction, this is done precisely with those contacts that are usually named as ``friends'' and which are seen outside the home and the neighborhood \citep{dimaggio87, fiske87, frith98}.  However, there is nothing inherent in the logic of the network theory of taste transmission that makes it incompatible with current data suggesting large rates of network turnover and the liability of newness of social ties.  

The empirical fact of network volatility coupled with the network assumption of taste as being largely a function of the cultural choices of the individual's current set of contacts, implies large-scale volatility of those very same cultural tastes.  Thus if we take the network perspective seriously, we should expect that as friends change, so do tastes, so we should expect small to moderate correlations between an individual's taste at time t and her taste at time $t+n$. The reason for this is that as \citet{mark03} notes a taste that is not ``reinforced'' by social network influence through interaction is likely to be ``forgotten'' (e.g. lost).   This is a key assumption of the network based model of taste acquisition.  Thus, if we retain the network assumption that individual tastes are largely a function of the tastes displayed by the persons current network alters, and if we accept the fact that these alters are constantly change over the adult life course, then we should expect cultural tastes to themselves exhibit high volatility over time.